%-------------------MISC-------------------
 
% Langkah 1: Buat counter baru untuk label otomatis
\newcounter{autofigure}
\newcounter{autotable}

% Mengatur agar caption di slide menampilkan nomor
\setbeamertemplate{caption}[numbered]
\setbeamertemplate{caption label separator}{: }

\makeatletter
\newcommand*{\rom}[1]{\expandafter\@slowromancap\romannumeral #1@}
\makeatother
%\setbeamertemplate{caption}[numbered] %Numbering caption
\renewcommand{\figurename}{Gambar} %Change "Figure" in caption to "Gambar"

\setbeamertemplate{bibliography item}{\insertbiblabel}
\renewcommand*{\bibfont}{\scriptsize}
\setbeamercolor{bibliography entry author}{fg=black}
\setbeamercolor{bibliography entry title}{fg=black} 
\setbeamercolor{bibliography entry location}{fg=black} 
\setbeamercolor{bibliography entry note}{fg=black}  
 
%----------------Styling----------------
\makeatletter
% --- Tabel ---%
\addto\captionsindonesian{ \renewcommand{\tablename}{Tabel} }
\newenvironment{Tabel}[1]{
  \begin{table}[!htbp]
  \centering
  \refstepcounter{autotable}
  \hypertarget{target.tab.\theautotable}{}
  \caption{#1}\label{tab.auto.\theautotable}
  % Menambahkan entri ke .lot dengan perintah hyperlink eksplisit
  \addcontentsline{lot}{table}{\protect\numberline{Tabel \thetable: \protect\hyperlink{target.tab.\thetable}{#1}}}
  \begin{tabular}{ll}
  \toprule
}{
  \bottomrule
  \end{tabular}
  \end{table}
}

% Memaksa \listoffigures & \listoftables untuk membaca file yang benar
\renewcommand{\listoffigures}{\@starttoc{lof}}
\renewcommand{\listoftables}{\@starttoc{lot}}
\makeatother

% --- Definisi, Aksioma, Teorema, dan Bukti --- %
\newcounter{nomorTeorema}
\newcounter{nomorAksioma} 
\newcounter{nomorDefinisi}

\NewDocumentEnvironment{Teorema}{ O{} O{} }{\only<1>{
      \stepcounter{nomorTeorema}
      \ifstrempty{#2}{}{\label{#2}
      }
    } 
    \begin{tcolorbox}[enhanced, colback=teal!10, colframe=teal!80!black, fonttitle=\bfseries, coltitle=white, colbacktitle=teal!70!black, left=3mm, right=3mm, top=2mm, bottom=2mm, title={Teorema \thenomorTeorema\ifstrempty{#1}{}{:\ #1}}] \setbeamertemplate{enumerate items}[default] \setlength{\leftmargini}{6mm}
}{  \end{tcolorbox}}

\newenvironment{Aksioma}[1][]{ \only<1>{\stepcounter{nomorAksioma}} \begin{tcolorbox}[ enhanced, colback=purple!10, colframe=purple!80!black, fonttitle=\bfseries, coltitle=white, colbacktitle=purple!70!black, left=3mm, right=3mm, top=2mm, bottom=2mm, title={Aksioma \thenomorAksioma\ifstrempty{#1}{}{:\ #1}} ] \setbeamertemplate{enumerate items}[default]    \setlength{\leftmargini}{6mm} 
}{ \end{tcolorbox} }

\newenvironment{Definisi}[1][]{ \only<1>{\stepcounter{nomorDefinisi}} \begin{tcolorbox}[ enhanced, colback=orange!10, colframe=orange!80!black, fonttitle=\bfseries, coltitle=white, colbacktitle=orange!70!black, left=3mm, right=3mm, top=2mm, bottom=2mm, title={Definisi \thenomorDefinisi\ifstrempty{#1}{}{:\ #1}} ] \setbeamertemplate{enumerate items}[default]    \setlength{\leftmargini}{6mm} 
}{ \end{tcolorbox} }

\usetikzlibrary{trees, shapes.geometric, shapes.symbols, shapes.misc, arrows.meta, positioning, matrix, backgrounds, calc}

\usepackage{pgfplots}
% Optional but highly recommended for stable results
\pgfplotsset{compat=newest} 
\usepgfplotslibrary{fillbetween}

\newcommand{\drawDogearNote}[3]{%
    \node (#2) [text width=2.5cm, align=left, font=\scriptsize] at #1 {#3};
    \def\foldsize{0.4cm}
    \begin{scope}[on background layer]
        \draw[draw, fill=yellow!40] 
            ($(#2.south west)+(-0.2,-0.2)$) -- ($(#2.south east)+(0.2,-0.2)$) -- 
            ($(#2.north east)+(0.2,0.2-\foldsize)$) -- ($(#2.north east)+(-\foldsize+0.2,0.2)$) -- 
            ($(#2.north west)+(-0.2,0.2)$) -- cycle;
        \draw ($(#2.north east)+(0.2,0.2-\foldsize)$) -- 
              ($(#2.north east)+(-\foldsize+0.2,0.2-\foldsize)$) -- 
              ($(#2.north east)+(-\foldsize+0.2,0.2)$);
    \end{scope}
}

\tikzset{
    startstop/.style = {ellipse, rounded corners, minimum width=3cm, minimum height=1cm, text centered, draw=black},
    process/.style = {rectangle, minimum width=3cm, minimum height=1cm, text centered, draw=black},
    decision/.style  = {diamond, aspect=2, minimum width=3cm, minimum height=1cm, text centered, draw=black},
    io/.style = {trapezium, trapezium left angle=70, trapezium right angle=110, minimum width=3cm, minimum height=1cm, text centered, draw=black},
    connector/.style = {circle, minimum size=1cm, text centered, draw=black},
    document/.style  = {
        tape, 
        tape bend top=none, % Hanya bagian bawah yang melengkung
        draw, 
        minimum width=3cm, 
        minimum height=2cm, 
        text centered
    }, 
    database/.style = {cylinder, shape border rotate=90, aspect=0.25, minimum height=1.5cm, minimum width=2.5cm, text centered, draw=black},
    arrow/.style = {thick, ->, >=stealth}, 
    action/.style = {rectangle, rounded corners=0.5cm, draw, fill=orange!30, minimum width=2.5cm, minimum height=1cm, text centered},
    activity/.style = {rectangle, rounded corners=0.2cm, draw, fill=blue!20, minimum width=3cm, minimum height=1.2cm, text centered},
    initial/.style = {circle, draw, radius=5pt, fill=black, inner sep=2pt},
    final/.style = {circle, draw=white, fill=black, double=white, double=black, double distance=1pt, inner sep=2.5pt},
    fork_join_v/.style={rectangle, draw, fill=black, minimum width=1.5mm, minimum height=3cm},     fork_join_h/.style={rectangle, draw, fill=black, minimum height=1.5mm, minimum width=3cm}, 
    decision/.style  = {diamond, aspect=2, draw, fill=green!30, text centered},
    object_node/.style = {rectangle, draw, fill=blue!20, text centered},
    note/.style = {note, draw, fill=yellow!40, text width=2.5cm},     
    lane_text/.style = {font=\bfseries}    
}

\makeatletter
\newcommand{\listofequations}{\@starttoc{loe}}
\newcommand{\addequation}[1]{
    \addcontentsline{loe}{section}{\protect\numberline{\theequation}#1}
}
\makeatother

\newcommand{\slideQA}{
\begin{frame}{Q\&A}
    \begin{center}
        \Huge Ada Pertanyaan?
    \end{center}
\end{frame}
}